%% 2-11 なぜ文字を使うの？　第3部
%% ことばの式 → 表 → □ → 文字。計算の順番のきまりと、そのわけ（第11回 X12）。いつでも成り立つきまり（交換法則・分配法則）。
%% この章ではまだ × を省かない（× を省くやくそくは 2-12）。
\makeatletter
% つまようじの色
\colorlet{bkiilwood}{bkStem!45!bkYellow}
% つまようじ1本  \bkiilpick{(x1,y1)}{(x2,y2)}{色}   両はしを少しあけて、1本ずつに見えるようにする
\newcommand\bkiilpick[3]{\draw[#3, line width=1.7pt, line cap=round, shorten <=0.75mm, shorten >=0.75mm] #1 -- #2;}
% つながった正方形  \bkiilsq{(x,y)}{正方形の数}{1辺mm}{はじめの1本の色}{3本ずつの色}
\newcommand\bkiilsq[5]{\begin{scope}[shift={#1}]
  \bkiilpick{(0,0)}{(0,#3)}{#4}
  \foreach \bkiil@k in {1,...,#2} {
    \bkiilpick{({(\bkiil@k-1)*#3},#3)}{({\bkiil@k*#3},#3)}{#5}
    \bkiilpick{({(\bkiil@k-1)*#3},0)}{({\bkiil@k*#3},0)}{#5}
    \bkiilpick{({\bkiil@k*#3},0)}{({\bkiil@k*#3},#3)}{#5}}
\end{scope}}
% はなれた正方形（4本ずつ）  \bkiilsqsep{(x,y)}{正方形の数}{1辺mm}{間mm}{色}
\newcommand\bkiilsqsep[5]{\begin{scope}[shift={#1}]
  \foreach \bkiil@k in {1,...,#2} {\begin{scope}[shift={({(\bkiil@k-1)*(#3+#4)},0)}]
    \bkiilpick{(0,0)}{(0,#3)}{#5}\bkiilpick{(0,#3)}{(#3,#3)}{#5}\bkiilpick{(0,0)}{(#3,0)}{#5}\bkiilpick{(#3,0)}{(#3,#3)}{#5}
  \end{scope}}
\end{scope}}
% 自転車（横から）  \bkiilbike{(x,y)}{大きさ（1で長さ約24mm）}{色}   (x,y) は後ろのタイヤの下
\newcommand\bkiilbike[3]{\begin{scope}[shift={#1}, scale=#2]
  \draw[bkInk!80, line width=1.1pt] (4.5,4.5) circle (4.2) (19.5,4.5) circle (4.2);
  \draw[#3, line width=1.4pt, line join=round] (4.5,4.5) -- (10.5,4.5) -- (8,11) -- cycle (10.5,4.5) -- (17.5,11) -- (8,11);
  \draw[#3, line width=1.4pt] (17.5,11) -- (19.5,4.5) (17.5,11) -- (16.6,13.6);
  \draw[bkInk!80, line width=1.3pt, line cap=round] (15.2,13.8) -- (18.2,13.4);
  \draw[bkInk!80, line width=1.6pt, line cap=round] (6.4,12.4) -- (9.6,12.4);
  \draw[#3, line width=1.2pt] (8,11) -- (8,12.2);
\end{scope}}
% みかん  \bkiilorange{(x,y)}{半径}
\newcommand\bkiilorange[2]{\begin{scope}[shift={#1}]
  \filldraw[fill=bkOrange, draw=bkOrange!70!black, line width=0.5pt] (0,0) ellipse ({#2} and {#2*0.85});
  \fill[bkLeaf] (0,{#2*0.85}) ellipse ({#2*0.35} and {#2*0.16});
\end{scope}}
% ケーキ（ショートケーキ1切れ）  \bkiilcake{(x,y)}{はば}   (x,y) は左下
\newcommand\bkiilcake[2]{\begin{scope}[shift={#1}, scale=#2/10]
  \filldraw[fill=bkYellow!35, draw=bkStem!70, line width=0.5pt] (0,0) -- (10,0) -- (10,5) -- (0,7.6) -- cycle;
  \fill[white] (0,7.6) -- (10,5) -- (10,5.9) -- (0,8.5) -- cycle;
  \draw[bkStem!70, line width=0.5pt] (0,8.5) -- (10,5.9) -- (10,0);
  \fill[white] (0,3.4) -- (10,2.2) -- (10,3.2) -- (0,4.4) -- cycle;
  \fill[bkRed] (2.4,9.6) circle (1.3);
\end{scope}}
% ことばの式を、まん中に1行で置く  \bkiilline{…}
\newcommand\bkiilline[1]{\par\addvspace{3pt}{\centering #1\par}\addvspace{3pt}}
% 点（おはじき）を、たて a こ × よこ b れつ に並べる  \bkiildots{(x,y)}{たて}{よこ}{色}
\newcommand\bkiildots[4]{\begin{scope}[shift={#1}]
  \foreach \bkiil@i in {1,...,#3} {\foreach \bkiil@j in {1,...,#2} {
    \filldraw[fill=#4, draw=bkInk!45, line width=0.4pt] ({(\bkiil@i-1)*4.6},{(\bkiil@j-1)*4.6}) circle (1.7);}}
\end{scope}}
\makeatother

\bkchapter{2}{11}{なぜ\R{文字}{も|じ}を\R{使}{つか}うの？}{\bklevel{2}}{なぜ\R{文字}{も|じ}を\R{使}{つか}うの？}
\bkmichishirube{1-12,1-13,2-2,2-3}{2-12,2-13,2-16,3-3}

\begin{goal}
\begin{itemize}[label={\color{bkTheme}●}, leftmargin=1.6\zw, itemsep=0pt, topsep=1pt]
\item \bklevel{2}\enspace ことばの\R{式}{しき}を、□や\R{文字}{も|じ}を\R{使}{つか}った\R{式}{しき}で\R{表}{あらわ}せる。
\item \bklevel{1}\enspace \R{計算}{けい|さん}の\R{順番}{じゅん|ばん}のきまりと、そのわけがわかり、1\R{行}{ぎょう}の\R{式}{しき}を\R{正}{ただ}しく\R{計算}{けい|さん}できる。
\item \bklevel{2}\enspace どんな\R{数}{かず}でも\R{成}{な}り\R{立}{た}つきまりを、\R{文字}{も|じ}を\R{使}{つか}って1\R{行}{ぎょう}で\R{書}{か}ける。
\end{itemize}
\end{goal}

\begin{hatena}
\R{歩}{ある}く\R{人}{ひと}・\R{自転車}{じ|てん|しゃ}・\R{自動車}{じ|どう|しゃ}が、それぞれ\R{一定}{いっ|てい}の\R{速}{はや}さで\R{進}{すす}んでいます。\R{進}{すす}んだ\R{道}{みち}のりを\R{表}{ひょう}にすると、\R{次}{つぎ}のようになります。
\begin{bkcenter}
\begin{bktable}{colspec={Q[l,m,wd=43mm] *{4}{Q[c,m,wd=12mm]} Q[c,m,wd=8mm]}, rowsep=1.5pt, colsep=4.5pt,
  row{1}={bg=bkYellow!30, font=\sffamily\bfseries\small}, cells={font=\small}}
 & 1\R{秒}{びょう} & 2\R{秒}{びょう} & 3\R{秒}{びょう} & 4\R{秒}{びょう} & … \\
\tikz[baseline=1.2mm, x=1mm, y=1mm]{\bkkid{(4,0)}{8.5}{bkGreen}{bkHair}\path (0,0) (17,0);}\R{歩}{ある}く\R{人}{ひと}（\qty{1}{m/s}） & \qty{1}{m} & \qty{2}{m} & \qty{3}{m} & \qty{4}{m} & … \\
\tikz[baseline=1.2mm, x=1mm, y=1mm]{\bkiilbike{(0,0)}{0.62}{bkBlue}\path (17,0);}\R{自転車}{じ|てん|しゃ}（\qty{5}{m/s}） & \qty{5}{m} & \qty{10}{m} & \qty{15}{m} & \qty{20}{m} & … \\
\tikz[baseline=1.2mm, x=1mm, y=1mm]{\bkcar{(1,0)}{0.5}{bkRed}\path (17,0);}\R{自動車}{じ|どう|しゃ}（\qty{15}{m/s}） & \qty{15}{m} & \qty{30}{m} & \qty{45}{m} & \qty{60}{m} & … \\
\end{bktable}
\end{bkcenter}
\R{速}{はや}さも\R{時間}{じ|かん}も、ほかにいくらでもあるので、\R{表}{ひょう}はどこまでも\R{大}{おお}きくなります。この\R{表}{ひょう}に\R{入}{はい}る\R{道}{みち}のりを\R{全部}{ぜん|ぶ}、たった1\R{行}{ぎょう}の\R{式}{しき}で\R{書}{か}けるかな？
\end{hatena}

\R{文字}{も|じ}を\R{使}{つか}うと、それができます。\R{計算}{けい|さん}の\R{順番}{じゅん|ばん}のきまりと、そのわけも\R{学}{まな}びます。

%% ---------------- やってみよう・思い出そう ----------------
\begin{yattemiyou}[つまようじで\R{正方形}{せい|ほう|けい}をつなげよう]
\begin{bkfigure}
\begin{tikzpicture}[x=1mm, y=1mm]
  \bkiilsq{(0,0)}{1}{14}{bkiilwood}{bkiilwood}
  \bkiilsq{(29,0)}{2}{14}{bkiilwood}{bkiilwood}
  \bkiilsq{(72,0)}{3}{14}{bkiilwood}{bkiilwood}
  \node[bk label, anchor=north] at (7,-2.5) {1こ};
  \node[bk label, anchor=north] at (43,-2.5) {2こ};
  \node[bk label, anchor=north] at (93,-2.5) {3こ};
  \node[bk label, anchor=west] at (117,7) {…};
\end{tikzpicture}
\end{bkfigure}
\begin{enumerate}[label={\bkcircled[bkGreen]{\arabic*}}, leftmargin=2.4\zw, itemsep=1pt, topsep=2pt]
\item つまようじで、\R{正方形}{せい|ほう|けい}を\R{横}{よこ}に1こ、2こ、3こ…とつなげてならべよう。
\item \R{正方形}{せい|ほう|けい}の\R{数}{かず}と、\R{使}{つか}ったつまようじの\R{本数}{ほん|すう}を、ノートの\R{表}{ひょう}に\R{書}{か}きこもう。
\item \R{正方形}{せい|ほう|けい}を1こふやすと、つまようじは\R{何}{なん}\R{本}{ぼん}ふえるかな？
\item \R{正方形}{せい|ほう|けい}を100こつなげるには、つまようじが\R{何}{なん}\R{本}{ぼん}いるかな？ならべずに、\R{計算}{けい|さん}で\R{考}{かんが}えよう。
\end{enumerate}
\begin{bkcenter}
\begin{bktable}{colspec={Q[l,m,wd=40mm] *{5}{Q[c,m,wd=9mm]} Q[c,m,wd=12mm]}, rowsep=1.5pt, colsep=4pt,
  row{1}={bg=bkGreen!14, font=\sffamily\bfseries\small}, cells={font=\small}}
\R{正方形}{せい|ほう|けい}の\R{数}{かず}（こ） & 1 & 2 & 3 & 4 & 5 & 100 \\
つまようじの\R{本数}{ほん|すう}（\R{本}{ほん}） & & & & & & \\
\end{bktable}
\end{bkcenter}
\end{yattemiyou}

\begin{omoidasou}
\bkrecall{1-12}{かけ\R{算}{ざん}は「1つ\R{分}{ぶん}の\R{数}{かず}×いくつ\R{分}{ぶん}」。$3\times4$ は「3が4つ\R{分}{ぶん}」の、\R{同}{おな}じ\R{数}{かず}のまとまり。}
\bkrecall{1-13}{\mbox{$3\times□=12$} の□に\R{入}{はい}る\R{数}{かず}は、\R{九九}{く|く}の\R{表}{ひょう}でさがせる（□は4）。}
\bkrecall{2-3}{\mbox{\R{道}{みち}のり＝\R{速}{はや}さ×\R{時間}{じ|かん}}。\qty{5}{m/s}で4\R{秒}{びょう}\R{進}{すす}むと、$\qty{5}{m/s}\times\qty{4}{s}=\qty{20}{m}$。}
\bkrecall{2-2}{マイナスの\R{数}{かず}のかけ\R{算}{ざん}：$(-2)\times3=-6$、$(-2)\times(-3)=6$。ひき\R{算}{ざん}は、マイナスの\R{数}{かず}をたすのと\R{同}{おな}じ：$10-3=10+(-3)$}
\end{omoidasou}

\bkfukidashi{mina}{100こなんて、ならべるのは\R{大変}{たい|へん}！　\R{表}{ひょう}をながめていると、きまりが\R{見}{み}えてこないかな？}

\clearpage
%% ---------------- 数学の道具 ----------------
\bksection{doguu}{\R{数学}{すう|がく}の\R{道具}{どう|ぐ}}

\bksubsection{1}{ことばの\R{式}{しき}を、□で\R{短}{みじか}く\R{書}{か}く}

\leavevmode\label{bk:2-11:kotoba}\bkhatenaref の\R{表}{ひょう}の\R{数}{かず}は、どれも\R{同}{おな}じ\R{計算}{けい|さん}で\R{出}{だ}せます。\R{自転車}{じ|てん|しゃ}の3\R{秒}{びょう}は $5\times3=15$、\R{自動車}{じ|どう|しゃ}の4\R{秒}{びょう}は $15\times4=60$。ことばで\R{書}{か}くと、
\bkiilline{\R{道}{みち}のり＝\R{速}{はや}さ×\R{時間}{じ|かん}}
です（2-3）。このように、ことばで\R{書}{か}いた\R{式}{しき}を\textbf{ことばの\R{式}{しき}}といいます。\R{意味}{い|み}はよくわかりますが、\R{書}{か}くのに\R{時間}{じ|かん}がかかります。そこで、ことばのかわりに□・○・△を\R{使}{つか}ってみましょう。

\noindent\begin{minipage}[c]{0.56\linewidth}
\bkiilline{\large ○＝□×△}
（○は\R{道}{みち}のり、□は\R{速}{はや}さ、△は\R{時間}{じ|かん}）\par\vspace{2pt}
□は「\R{数}{かず}を\R{入}{い}れる\R{箱}{はこ}」です。□に5、△に4を\R{入}{い}れると、\mbox{○＝5×4＝20}。\R{自転車}{じ|てん|しゃ}の4\R{秒}{びょう}と\R{同}{おな}じです。
\end{minipage}\hfill
\begin{minipage}[c]{0.4\linewidth}\centering
\begin{tikzpicture}[x=1mm, y=1mm]
  \foreach \v/\xx in {1/4,5/19,15/34} {\node[draw=bkGray, line width=0.7pt, fill=white, rounded corners=1mm, minimum width=9mm, minimum height=6.5mm, font=\sffamily\bfseries\small] (c\v) at (\xx,19) {\v};}
  \node[draw=bkTheme, line width=1.4pt, minimum width=12mm, minimum height=10mm, fill=bkTheme!8] (box) at (19,0) {};
  \foreach \v in {1,5,15} \draw[bk link] (c\v) -- (box);
\end{tikzpicture}
\bkfigcaption{□には、どんな\R{速}{はや}さでも\R{入}{はい}る。}
\end{minipage}

\R{第}{だい}1\R{巻}{かん}の「つぎのステージへ」のカードにあった \mbox{□＋3＝8} の□も、\R{数}{かず}を\R{入}{い}れる\R{箱}{はこ}です。□に5を\R{入}{い}れると $5+3=8$ でぴったり\R{合}{あ}うので、□は5です。

\bksubsection{2}{□のかわりに、\R{文字}{も|じ}を\R{使}{つか}う}

\leavevmode\label{bk:2-11:moji}□・○・△がたくさん\R{出}{で}てくると、どれが\R{何}{なん}の\R{箱}{はこ}だったか、わからなくなります。そこで、\R{数学}{すう|がく}や\R{物理}{ぶつ|り}では、□のかわりにアルファベットの\R{文字}{も|じ}を\R{使}{つか}います。\R{物理}{ぶつ|り}では、\R{量}{りょう}を\R{表}{あらわ}す\R{英語}{えい|ご}の\R{頭文字}{かしら|も|じ}をよく\R{使}{つか}います。2-3で\R{量}{りょう}の\R{名前}{な|まえ}として\R{使}{つか}った $x$、$t$ を、ここからは\R{数}{かず}の\R{代}{か}わりの\R{文字}{も|じ}として\R{計算}{けい|さん}に\R{使}{つか}います。

\begin{bkcenter}
\begin{bktable}{colspec={Q[l,m,wd=19mm] *{2}{Q[c,m,wd=17mm]} Q[c,m,wd=23mm] *{2}{Q[c,m,wd=17mm]}}, rowsep=1pt, colsep=5pt, cells={font=\small}, row{1}={font=\large, bg=bkTheme!12}, row{3}={font=\footnotesize}}
\R{文字}{も|じ} & $t$ & $v$ & $x$ & $m$ & $F$ \\
\R{表}{あらわ}す\R{量}{りょう} & \R{時間}{じ|かん} & \R{速}{はや}さ & \R{位置}{い|ち}・\R{道}{みち}のり & \R{質量}{しつ|りょう} & \R{力}{ちから}（2-15） \\
もとの\R{英語}{えい|ご} & \textbf{t}ime & \textbf{v}elocity & （\R{座標}{ざ|ひょう}の\R{文字}{も|じ}） & \textbf{m}ass & \textbf{F}orce \\
\end{bktable}
\bktabcaption{\R{物理}{ぶつ|り}でよく\R{使}{つか}う\R{文字}{も|じ}}
\end{bkcenter}

\R{単位}{たん|い}の m（メートル）はまっすぐな\R{字}{じ}、\R{量}{りょう}の $m$（\R{質量}{しつ|りょう}）はななめの\R{字}{じ}で\R{書}{か}いて、\R{区別}{く|べつ}します。この\R{文字}{も|じ}を\R{使}{つか}うと、ことばの\R{式}{しき}は、\R{次}{つぎ}のように\R{短}{みじか}く\R{書}{か}けます。
\begin{bkfigure}
\begin{tikzpicture}[x=1mm, y=1mm]
  \foreach \y/\a/\b/\c/\n in {16/\R{道}{みち}のり/\R{速}{はや}さ/\R{時間}{じ|かん}/ことばの\R{式}{しき}, 8/○/□/△/\R{記号}{き|ごう}の\R{式}{しき}} {
    \node[anchor=east, font=\sffamily\small, text=bkGray] at (14,\y) {\n};
    \node[font=\small] at (30,\y) {\a};
    \node[font=\small] at (42,\y) {＝};
    \node[font=\small] at (54,\y) {\b};
    \node[font=\small] at (66,\y) {×};
    \node[font=\small] at (78,\y) {\c};}
  \fill[bkTheme!10, rounded corners=1.5mm] (19,-3.8) rectangle (89,3.8);
  \node[anchor=east, font=\sffamily\bfseries\small, text=bkTheme] at (14,0) {\R{文字}{も|じ}の\R{式}{しき}};
  \node[font=\large] at (30,0) {$x$}; \node[font=\large] at (42,0) {$=$}; \node[font=\large] at (54,0) {$v$};
  \node[font=\large] at (66,0) {$\times$}; \node[font=\large] at (78,0) {$t$};
\end{tikzpicture}
\bkfigcaption{\R{同}{おな}じことを、3つの\R{書}{か}き\R{方}{かた}で\R{書}{か}いた。\R{下}{した}へいくほど\R{短}{みじか}い。}
\end{bkfigure}

\bkhatenaref の\R{答}{こた}えは、この1\R{行}{ぎょう} \mbox{$x=v\times t$} です。$v$ に1を\R{入}{い}れれば\R{歩}{ある}く\R{人}{ひと}の\R{行}{ぎょう}、15を\R{入}{い}れれば\R{自動車}{じ|どう|しゃ}の\R{行}{ぎょう}。\R{表}{ひょう}の\R{全部}{ぜん|ぶ}が、この1\R{行}{ぎょう}に\R{入}{はい}っています。

\R{文字}{も|じ}は、\textbf{どんな\R{数}{かず}でも\R{入}{はい}る\R{箱}{はこ}}です。このように、\R{文字}{も|じ}を\R{使}{つか}って\R{書}{か}いた\R{式}{しき}を\textbf{\mbox{\R{文字式}{も|じ|しき}}}といいます。\bkidx{もじしき}{\R{文字式}{も|じ|しき}}

\bksubsection{3}{\R{計算}{けい|さん}の\R{順番}{じゅん|ばん}：きまりがないと、\R{答}{こた}えが2つになる}

\leavevmode\label{bk:2-11:junban}\bkyattemiyouref のつまようじも、\R{文字}{も|じ}で\R{表}{あらわ}してみましょう。\R{正方形}{せい|ほう|けい}の\R{数}{かず}を $n$（\textbf{n}umber の\R{頭文字}{かしら|も|じ}）とします。\R{下}{した}の\R{図}{ず}のAのように、\R{左}{ひだり}はしの1\R{本}{ぽん}に、\R{正方形}{せい|ほう|けい}1こにつき3\R{本}{ぼん}ずつがつながるので、
\bkiilline{\R{本数}{ほん|すう}＝$1+3\times n$\qquad（$n=1$ で4、$n=2$ で7、$n=3$ で10）}
$n=100$ なら $1+3\times100$ です。これを\R{左}{ひだり}から\R{順}{じゅん}に\R{計算}{けい|さん}すると、$1+3=4$、$4\times100=400$ で400\R{本}{ぽん}。ところが、\R{図}{ず}で\R{数}{かぞ}えると、\R{左}{ひだり}はしの1\R{本}{ぽん}と、3\R{本}{ぼん}ずつが100こで300\R{本}{ぽん}、\R{合}{あ}わせて301\R{本}{ぽん}です。どちらが\R{正}{ただ}しいかな？

\begin{bkfigure}
\begin{tikzpicture}[x=1mm, y=1mm]
  \node[anchor=west, font=\sffamily\bfseries\small, text=bkBlue!85!black] at (0,32) {A　つながった\R{正方形}{せい|ほう|けい}};
  \bkiilsq{(9,9)}{4}{11}{bkOrange}{bkBlue!70}
  \node[font=\sffamily\bfseries\footnotesize, text=bkOrange!85!black, anchor=east] at (7,14.5) {1\R{本}{ぽん}};
  \foreach \k in {1,2,3,4} \node[font=\sffamily\bfseries\footnotesize, text=bkBlue!85!black, anchor=south] at ({9+\k*11-5.5},21.5) {3\R{本}{ぼん}};
  \node[anchor=north west, font=\small] at (9,5) {$1+3\times n$　（$n=100$ で301\R{本}{ぽん}）};
  \node[anchor=west, font=\sffamily\bfseries\small, text=bkGreen!70!black] at (72,32) {B　はなれた\R{正方形}{せい|ほう|けい}};
  \bkiilsqsep{(72,9)}{4}{11}{3.5}{bkGreen!80!black}
  \node[anchor=north west, font=\small] at (72,5) {$(1+3)\times n$　（$n=100$ で400\R{本}{ぽん}）};
\end{tikzpicture}
\bkfigcaption{Aは $3\times n$ を\R{先}{さき}に\R{計算}{けい|さん}する。Bは $1+3$ を\R{先}{さき}に\R{計算}{けい|さん}する。}
\end{bkfigure}

\R{正}{ただ}しいのは301\R{本}{ぽん}。$3\times100$ を\R{先}{さき}に\R{計算}{けい|さん}しなければいけません。400\R{本}{ぽん}になるのは、Bのように、4\R{本}{ほん}ずつの\R{正方形}{せい|ほう|けい}をはなしてならべたときです。「$1+3$」をひとまとまりにして\R{先}{さき}に\R{計算}{けい|さん}したいときは、\textbf{かっこ}（　）でかこんで $(1+3)\times n$ と\R{書}{か}きます。

もう1つ\R{見}{み}てみましょう。$2+3\times5$ を、\R{2人}{ふたり}の\R{子}{こ}が\R{計算}{けい|さん}しました。
\begin{bkcenter}
\begin{bktable}{colspec={Q[l,m,wd=36mm] Q[l,m,wd=66mm]}, rowsep=2pt, cells={font=\small}}
ソラ（\R{左}{ひだり}から\R{順}{じゅん}に） & $2+3=5$、$5\times5=25$ \\
ミナ（×を\R{先}{さき}に） & $3\times5=15$、$2+15=17$ \\
\end{bktable}
\end{bkcenter}
\R{同}{おな}じ\R{式}{しき}なのに、\R{答}{こた}えが2つ\R{出}{で}てしまいました。これでは、\R{式}{しき}を\R{書}{か}いた\R{人}{ひと}が\R{何}{なに}を\R{伝}{つた}えたかったのか、わかりません。お\R{店}{みせ}のレジや、\R{橋}{はし}の\R{設計図}{せっ|けい|ず}の\R{計算}{けい|さん}で、\R{人}{ひと}によって\R{答}{こた}えがちがったら\R{大変}{たい|へん}です。そこで、\R{世界中}{せ|かい|じゅう}で\R{同}{おな}じ\R{順番}{じゅん|ばん}で\R{計算}{けい|さん}するやくそくを\R{決}{き}めました。\bkidx{けいさんのじゅんばん}{\R{計算}{けい|さん}の\R{順番}{じゅん|ばん}}

\begin{bkcenter}
\begin{tikzpicture}[x=1mm, y=1mm]
  \node[bk node, minimum width=34mm, minimum height=12mm] (a) at (0,0) {① かっこの\R{中}{なか}を\\\R{先}{さき}に\R{計算}{けい|さん}する};
  \node[bk node, minimum width=34mm, minimum height=12mm] (b) at (44,0) {② ×と÷を、\\＋と−より\R{先}{さき}に};
  \node[bk node, minimum width=34mm, minimum height=12mm] (c) at (88,0) {③ \R{同}{おな}じなかまは、\\\R{左}{ひだり}から\R{順}{じゅん}に};
  \draw[bk link] (a) -- (b); \draw[bk link] (b) -- (c);
\end{tikzpicture}
\end{bkcenter}

2-7の\R{累乗}{るい|じょう}（$3^2$ など）がある\R{式}{しき}では、\R{累乗}{るい|じょう}を①のかっこの\R{次}{つぎ}、②の×÷より\R{先}{さき}に\R{計算}{けい|さん}します（$2\times3^2=2\times9=18$。\R{章末}{しょう|まつ}のコラムも\R{見}{み}よう）。つまり、かっこ → \R{累乗}{るい|じょう} → ×÷ → ＋− の\R{順}{じゅん}です。

このやくそくでは、$2+3\times5=2+15=17$ で、ミナが\R{正}{ただ}しいことになります。でも、なぜ「×が\R{先}{さき}」なのでしょう。くじ\R{引}{び}きで\R{決}{き}めたのではありません。\R{次}{つぎ}のページから、わけを\R{見}{み}ていきます。

\bksubsection{4}{×と÷が\R{先}{さき}のわけ：かけ\R{算}{ざん}は1つのまとまり}

\leavevmode\label{bk:2-11:kakesaki}かけ\R{算}{ざん}は「1つ\R{分}{ぶん}の\R{数}{かず}×いくつ\R{分}{ぶん}」でした（1-12）。$100\times3$ は「100が3つ\R{分}{ぶん}」で、\R{同}{おな}じ\R{数}{かず}を\R{集}{あつ}めた\textbf{1つのまとまり}を\R{表}{あらわ}しています。3つの\R{場面}{ば|めん}で\R{確}{たし}かめましょう。

\textbf{\R{場面}{ば|めん}1　\R{買}{か}い\R{物}{もの}}\quad 1こ100\R{円}{えん}のりんごを3こと、1こ50\R{円}{えん}のみかんを1こ\R{買}{か}いました。\R{代金}{だい|きん}はいくらかな？
\begin{bkfigure}
\begin{tikzpicture}[x=1mm, y=1mm]
  \foreach \k in {0,1,2} {\bkapple{(\k*14+6,6)}{5}{bkRed}
    \node[anchor=north, font=\sffamily\footnotesize, inner sep=0pt] at (\k*14+6,-1.5) {100\R{円}{えん}};}
  \bkiilorange{(66,5.5)}{4.6}
  \node[anchor=north, font=\sffamily\footnotesize, inner sep=0pt] at (66,-1.5) {50\R{円}{えん}};
  \draw[bkRed!80!black, line width=0.8pt, dashed, rounded corners=2mm] (-2,-6.5) rectangle (42,13.5);
  \node[anchor=south, font=\sffamily\bfseries\footnotesize, text=bkRed!80!black, inner sep=0pt] at (20,15.2) {$100\times3=300$　りんご\R{全部}{ぜん|ぶ}の\R{代金}{だい|きん}（1つのまとまり）};
  \node[anchor=west, font=\small, inner sep=0pt] at (78,5) {$100\times3+50$};
  \node[anchor=west, font=\small, inner sep=0pt] at (78,-1) {$=300+50=350$};
\end{tikzpicture}
\bkfigcaption{$100\times3$ は「りんご\R{全部}{ぜん|ぶ}の\R{代金}{だい|きん}」という1つの\R{量}{りょう}。それにみかんの\R{代金}{だい|きん}をたす。}
\end{bkfigure}
$100\times3$ は「りんご\R{全部}{ぜん|ぶ}の\R{代金}{だい|きん}」という1つの\R{量}{りょう}です。だから\R{先}{さき}に300という1つの\R{数}{かず}にまとめて、そこにみかんの50\R{円}{えん}をたします。みかんを\R{先}{さき}に\R{書}{か}いて $50+100\times3$ としても、\R{代金}{だい|きん}は\R{同}{おな}じ350\R{円}{えん}のはずです。もし\R{左}{ひだり}から\R{計算}{けい|さん}すると $50+100=150$、$150\times3=450$ になります。450\R{円}{えん}は「みかんとりんごを1こずつ\R{組}{くみ}にして、3\R{組}{くみ}\R{買}{か}った」ときの\R{代金}{だい|きん}で、\R{買}{か}った\R{物}{もの}とちがってしまいます。

\textbf{\R{場面}{ば|めん}2　L\R{字}{じ}の\R{部屋}{へ|や}の\R{面積}{めん|せき}}\quad \R{下}{した}の\R{図}{ず}の\R{部屋}{へ|や}の\R{面積}{めん|せき}は\R{何}{なん}\unit{m^2}かな？　2つの\R{長方形}{ちょう|ほう|けい}に\R{分}{わ}けて\R{考}{かんが}えます。

\noindent\begin{minipage}[c]{0.42\linewidth}\centering
\begin{tikzpicture}[x=1mm, y=1mm]
  \fill[bkBlue!20] (0,0) rectangle (24,18);
  \fill[bkOrange!25] (0,18) rectangle (12,30);
  \foreach \x in {6,12,18} \draw[bkLine, line width=0.4pt] (\x,0) -- (\x,18);
  \foreach \y in {6,12} \draw[bkLine, line width=0.4pt] (0,\y) -- (24,\y);
  \draw[bkLine, line width=0.4pt] (6,18) -- (6,30) (0,24) -- (12,24);
  \draw[bkInk, line width=0.9pt] (0,0) -- (24,0) -- (24,18) -- (12,18) -- (12,30) -- (0,30) -- cycle;
  \draw[bkInk, line width=0.6pt, dashed] (0,18) -- (12,18);
  \node[font=\sffamily\bfseries\small, text=bkBlue!80!black, fill=bkBlue!20, inner sep=1pt] at (12,9) {$4\times3$};
  \node[font=\sffamily\bfseries\small, text=bkOrange!80!black, fill=bkOrange!25, inner sep=1pt] at (6,24) {$2\times2$};
  \node[anchor=north, bk label] at (12,-1.5) {\qty{4}{m}};
  \node[anchor=west, bk label] at (25.5,9) {\qty{3}{m}};
  \node[anchor=south, bk label] at (6,31.5) {\qty{2}{m}};
  \node[anchor=west, bk label] at (13.5,24) {\qty{2}{m}};
\end{tikzpicture}
\bkfigcaption{1マスは\qty{1}{m^2}。}
\end{minipage}\hfill
\begin{minipage}[c]{0.55\linewidth}
\R{下}{した}の\R{長方形}{ちょう|ほう|けい}は $4\times3$、\R{上}{うえ}の\R{正方形}{せい|ほう|けい}は $2\times2$。どちらも、それだけで1つの\R{面積}{めん|せき}です。
\bkiilline{$4\times3+2\times2=12+4=16$}
マスを\R{数}{かぞ}えても16こなので、\qty{16}{m^2}。\par\vspace{2pt}
\R{左}{ひだり}から\R{順}{じゅん}に\R{計算}{けい|さん}すると、$4\times3=12$、$12+2=14$、$14\times2=28$ で\qty{28}{m^2}。\R{部屋}{へ|や}がすっぽり\R{入}{はい}る\R{横}{よこ}4\,m、たて5\,mの\R{長方形}{ちょう|ほう|けい}（\qty{20}{m^2}）より\R{広}{ひろ}くなり、おかしいとわかります。
\end{minipage}

\vspace{4pt}
\textbf{\R{場面}{ば|めん}3　\R{速}{はや}さと\R{時間}{じ|かん}}\quad \R{自転車}{じ|てん|しゃ}で\qty{5}{m/s}で60\R{秒}{びょう}\R{走}{はし}り、そのあと\R{自転車}{じ|てん|しゃ}をおりて\qty{1}{m/s}で100\R{秒}{びょう}\R{歩}{ある}きました。\R{全部}{ぜん|ぶ}で\R{何}{なん}m\R{進}{すす}んだかな？
\begin{bkfigure}
\begin{tikzpicture}[x=1mm, y=1mm]
  \draw[bkInk!70, line width=0.8pt] (0,0) -- (120,0);
  \fill[bkBlue!25] (0,0.5) rectangle (90,4);
  \fill[bkGreen!30] (90,0.5) rectangle (120,4);
  \bkiilbike{(30,4.5)}{0.75}{bkBlue}
  \bkkid{(105,4.5)}{13}{bkGreen}{bkHair}
  \node[anchor=north, font=\sffamily\footnotesize, inner sep=0pt] at (45,-1.8) {$5\times60=300$（m）};
  \node[anchor=north, font=\sffamily\footnotesize, inner sep=0pt] at (105,-1.8) {$1\times100=100$（m）};
  \draw[bkInk, line width=0.7pt] (90,-1) -- (90,5);
\end{tikzpicture}
\bkfigcaption{\R{自転車}{じ|てん|しゃ}で\R{進}{すす}んだ\R{道}{みち}のりと、\R{歩}{ある}いた\R{道}{みち}のりを、それぞれ\R{先}{さき}に\R{求}{もと}めてから\R{合}{あ}わせる。}
\end{bkfigure}
\R{道}{みち}のり＝\R{速}{はや}さ×\R{時間}{じ|かん}（2-3）なので、\R{自転車}{じ|てん|しゃ}の\R{道}{みち}のりと\R{歩}{ある}いた\R{道}{みち}のりを、それぞれ1つの\R{量}{りょう}として\R{先}{さき}に\R{計算}{けい|さん}します。
\bkiilline{$5\times60+1\times100=300+100=400$\qquad（\qty{400}{m}）}
\R{左}{ひだり}から\R{順}{じゅん}に\R{計算}{けい|さん}すると、$5\times60=300$、$300+1=301$、$301\times100=30100$ で\qty{30100}{m}（\R{約}{やく}\qty{30}{km}）。160\R{秒}{びょう}（3\R{分}{ぷん}たらず）で30\,kmも\R{進}{すす}むはずがありません。

\textbf{わり\R{算}{ざん}も\R{同}{おな}じ}\quad 600\R{円}{えん}でパンを3こ\R{買}{か}いました。1こ\R{分}{ぶん}の\R{値段}{ね|だん}は、600\R{円}{えん}を3つに\R{分}{わ}けた1つ\R{分}{ぶん}で、「1こあたりの\R{値段}{ね|だん}」という1つの\R{量}{りょう}です（2-3）。だから、パン1こと120\R{円}{えん}のジュースの\R{代金}{だい|きん}は、\R{先}{さき}に÷を\R{計算}{けい|さん}します。
\bkiilline{$600\div3+120=200+120=320$\qquad（320\R{円}{えん}）}

\begin{bkcenter}
{\setlength\fboxsep{4pt}\fcolorbox{bkTheme}{white}{\begin{minipage}{0.88\linewidth}
\par\addvspace{3pt}\textbf{×と÷が\R{先}{さき}のわけ}\par\noindent ×や÷でできた\R{部分}{ぶ|ぶん}（$100\times3$、$4\times3$、$5\times60$、$600\div3$）は、それぞれ「りんご\R{全部}{ぜん|ぶ}の\R{代金}{だい|きん}」「\R{長方形}{ちょう|ほう|けい}1つの\R{面積}{めん|せき}」のような\textbf{1つの\R{量}{りょう}}を\R{表}{あらわ}す。＋と−は、その\R{量}{りょう}どうしを\R{合}{あ}わせたり、ひいたりする。だから、まとまりを\R{先}{さき}に1つの\R{数}{かず}にしてから、＋と−をする。
\end{minipage}}}
\end{bkcenter}

\bksubsection{5}{かっこの\R{中}{なか}が\R{先}{さき}のわけ：かっこは「ひとまとまり」の\R{印}{しるし}}

\leavevmode\label{bk:2-11:kakko}\R{今度}{こん|ど}は、＋でできた\R{部分}{ぶ|ぶん}を、ひとまとまりとしてかけたいときを\R{考}{かんが}えます。おはじきを、たてに5こ、\R{横}{よこ}に5れつならべました。$5\times5=25$（こ）です。たての5こを、\R{青}{あお}2こと\R{赤}{あか}3こに\R{分}{わ}けて\R{見}{み}ると、1れつは $2+3$ こです。

\begin{bkfigure}
\begin{tikzpicture}[x=1mm, y=1mm]
  % 左：(2+3)×5
  \bkiildots{(2,2)}{2}{5}{bkBlue!70}
  \bkiildots{(2,11.2)}{3}{5}{bkRed!70}
  \draw[bkInk!60, line width=0.6pt, rounded corners=1mm] (-1.2,-1.2) rectangle (5.2,23.4);
  \node[anchor=north, font=\sffamily\footnotesize, inner sep=0pt] at (11.2,-3) {1れつ（1つ\R{分}{ぶん}）が $2+3$ こで、5れつ};
  \node[anchor=south, font=\small, inner sep=0pt] at (11.2,26.5) {$(2+3)\times5=5\times5=25$};
  \node[font=\sffamily\bfseries\small, text=bkGreen!60!black, inner sep=0pt] at (11.2,35) {○\ \R{正}{ただ}しい};
  % 右：2+3×5
  \begin{scope}[shift={(66,0)}]
  \bkiildots{(2,2)}{2}{1}{bkBlue!70}
  \bkiildots{(14,2)}{3}{5}{bkRed!70}
  \node[anchor=north, font=\sffamily\footnotesize, inner sep=0pt] at (20,-3) {\R{青}{あお}2こと、\R{赤}{あか}3こが5れつ};
  \node[anchor=south, font=\small, inner sep=0pt] at (20,26.5) {$2+3\times5=2+15=17$};
  \node[font=\sffamily\bfseries\small, text=bkRed, inner sep=0pt] at (20,35) {×\ 25にならない};
  \end{scope}
\end{tikzpicture}
\bkfigcaption{$2+3$ をひとまとまりにしてかけたいときは、かっこが\R{要}{い}る。}\label{bk:2-11:fig-kakko}
\end{bkfigure}

1れつ\R{分}{ぶん}の $2+3$ は、5という1つの\R{数}{かず}のかわりです。だから $5\times5$ は $(2+3)\times5$ と\R{書}{か}け、かっこの\R{中}{なか}を\R{先}{さき}に\R{計算}{けい|さん}して $5\times5=25$ になります。かっこをつけずに $2+3\times5$ と\R{書}{か}くと、②のきまりで $3\times5$ が\R{先}{さき}になり、17になります。25と17。\R{同}{おな}じおはじきの\R{数}{かず}のはずなのに、\R{答}{こた}えが\R{合}{あ}いません。$2+3\times5$ は、\R{図}{ず}\ref{bk:2-11:fig-kakko}の\R{右}{みぎ}がわのように「\R{青}{あお}2こと、\R{赤}{あか}3こずつ5れつ」という\R{別}{べつ}のならべ\R{方}{かた}を\R{表}{あらわ}しているからです。

\textbf{かっこ}は「ここをひとまとまりにして、\R{先}{さき}に\R{考}{かんが}えてね」という\R{印}{しるし}です。かっこの\R{中}{なか}は、1つの\R{数}{かず}のかわりなので、いちばん\R{先}{さき}に\R{計算}{けい|さん}して1つの\R{数}{かず}にします。

\textbf{ケーキとジュースの\R{代金}{だい|きん}}\quad 1\R{人分}{にん|ぶん}は、300\R{円}{えん}のケーキと100\R{円}{えん}のジュースです。2\R{人分}{にん|ぶん}の\R{代金}{だい|きん}はいくらかな？
\begin{bkfigure}
\begin{tikzpicture}[x=1mm, y=1mm]
  \foreach \k in {0,1} {\begin{scope}[shift={(\k*30,0)}]
    \bkiilcake{(1,0)}{11}
    \bkcup{(19,0)}{6.5}{10}{0.75}{0}
    \draw[bkTheme, line width=0.8pt, dashed, rounded corners=1.5mm] (-1.5,-2) rectangle (24.5,13.5);
    \node[anchor=north, font=\sffamily\footnotesize, inner sep=0pt] at (11.5,-3.5) {$300+100$};
  \end{scope}}
  \node[anchor=south, font=\small, inner sep=0pt] at (26.5,16) {$(300+100)\times2=400\times2=800$};
  \begin{scope}[shift={(72,0)}]
    \bkiilcake{(1,0)}{11}
    \bkcup{(19,0)}{6.5}{10}{0.75}{0}
    \bkcup{(29,0)}{6.5}{10}{0.75}{0}
    \node[anchor=north, font=\sffamily\footnotesize, inner sep=0pt] at (18,-3.5) {ケーキ1こと、ジュース2\R{本}{ほん}};
    \node[anchor=south, font=\small, inner sep=0pt] at (18,16) {$300+100\times2=500$};
  \end{scope}
\end{tikzpicture}
\bkfigcaption{\R{左}{ひだり}：1\R{人分}{にん|ぶん}をひとまとまりにして2つ\R{分}{ぶん}。\R{右}{みぎ}：かっこを\R{忘}{わす}れた\R{式}{しき}。}
\end{bkfigure}
「1\R{人分}{にん|ぶん}」の $300+100$ が、2つ\R{分}{ぶん}です。1\R{人分}{にん|ぶん}をひとまとまりにするので、かっこをつけて $(300+100)\times2=400\times2=800$（\R{円}{えん}）。かっこを\R{忘}{わす}れて $300+100\times2$ と\R{書}{か}くと500\R{円}{えん}になり、「ケーキ1こと、ジュース2\R{本}{ほん}」の\R{代金}{だい|きん}になってしまいます。

\begin{bkcenter}
{\setlength\fboxsep{4pt}\fcolorbox{bkTheme}{white}{\begin{minipage}{0.88\linewidth}
\par\addvspace{3pt}\textbf{かっこが\R{要}{い}るかどうかの\R{見分}{み|わ}け\R{方}{かた}}\par\noindent 「1\R{人分}{にん|ぶん}」「1れつ」「\R{全部}{ぜん|ぶ}で」のように、＋や−でできた\R{部分}{ぶ|ぶん}を\textbf{ひとまとまり}として、かけたり、わったり、ひいたりしたいときは、かっこでかこむ。
\end{minipage}}}
\end{bkcenter}

\bksubsection{6}{\R{同}{おな}じなかまは、\R{左}{ひだり}から\R{順}{じゅん}に}

\leavevmode\label{bk:2-11:hidari}\textbf{＋と−だけのとき}\quad 10\R{円}{えん}\R{持}{も}っていて、3\R{円}{えん}\R{使}{つか}い、そのあと2\R{円}{えん}もらいました。\R{今}{いま}いくら？
\begin{bkfigure}
\begin{tikzpicture}[x=1mm, y=1mm]
  \draw[bkInk, line width=0.7pt] (0,0) -- (118,0);
  \foreach \k in {0,...,12} {\draw[bkInk, line width=0.5pt] (\k*9,-1.2) -- (\k*9,1.2); \node[anchor=north, font=\sffamily\footnotesize, inner sep=0pt] at (\k*9,-2.2) {\k};}
  \draw[-{Stealth[length=2.2mm]}, bkRed, line width=1pt] (90,3) .. controls (85,9) and (68,9) .. (63,3);
  \node[anchor=south, font=\sffamily\bfseries\footnotesize, text=bkRed, inner sep=0pt] at (76.5,9) {$-3$};
  \draw[-{Stealth[length=2.2mm]}, bkBlue, line width=1pt] (63,-7.6) .. controls (66,-13) and (78,-13) .. (81,-7.6);
  \node[anchor=north, font=\sffamily\bfseries\footnotesize, text=bkBlue, inner sep=0pt] at (72,-13.6) {$+2$};
  \fill[bkTheme] (90,0) circle (1.3);
  \node[anchor=south, font=\sffamily\footnotesize, inner sep=0pt] at (98,2.6) {はじめ};
\end{tikzpicture}
\bkfigcaption{10から3もどって7、そこから2\R{進}{すす}んで9。}
\end{bkfigure}
お\R{金}{かね}は、\R{起}{お}きた\R{順}{じゅん}にふえたりへったりするので、$10-3+2$ は\R{左}{ひだり}から $10-3=7$、$7+2=9$ で9\R{円}{えん}。もし $3+2$ を\R{先}{さき}にすると $10-5=5$ になり、「5\R{円}{えん}\R{使}{つか}った」ことになってしまいます。「−3」の−は、すぐ\R{後}{あと}の3だけについた\R{印}{しるし}です。2-2のように $10+(-3)+2$ と\R{考}{かんが}えれば、どこから\R{計算}{けい|さん}しても9になります。

\textbf{×と÷だけのとき}\quad $24\div4\times2$ は「24を4つに\R{等}{ひと}しく\R{分}{わ}けた1つ\R{分}{ぶん}（6）の、2つ\R{分}{ぶん}」で、\R{左}{ひだり}から $24\div4=6$、$6\times2=12$。もし×を\R{先}{さき}にすると $24\div8=3$ で、「8つに\R{分}{わ}けた1つ\R{分}{ぶん}」という\R{別}{べつ}の\R{意味}{い|み}になります。8つに\R{分}{わ}けたいときは、かっこを\R{使}{つか}って $24\div(4\times2)$ と\R{書}{か}きます。

\begin{bkcenter}
\begin{tabular}{l@{\qquad}l}
$20-\underline{4\times3}=20-12=8$ & ②×を\R{先}{さき}に（\R{下線}{か|せん}の\R{所}{ところ}が\R{先}{さき}） \\[2pt]
$\underline{(20-4)}\times3=16\times3=48$ & ①かっこの\R{中}{なか}を\R{先}{さき}に \\[2pt]
$\underline{12\div4}\times5=3\times5=15$ & ③×と÷だけなので、\R{左}{ひだり}から \\[2pt]
$\underline{15-6}+4=9+4=13$ & ③＋と−だけなので、\R{左}{ひだり}から \\
\end{tabular}
\end{bkcenter}

\bksubsection{7}{\R{文字}{も|じ}で、いつでも\R{成}{な}り\R{立}{た}つきまりを\R{書}{か}く}

\leavevmode\label{bk:2-11:kimari}$2+3$ と $3+2$ は、どちらも5です。$4\times5$ と $5\times4$ も、どちらも20です（1-13）。ほかのどんな\R{数}{かず}でもそうなので、\R{文字}{も|じ}を\R{使}{つか}って \mbox{$a+b=b+a$}、\mbox{$a\times b=b\times a$} と1\R{行}{ぎょう}で\R{書}{か}けます。これを\textbf{\R{交換法則}{こう|かん|ほう|そく}}といいます。\bkidx{こうかんほうそく}{\R{交換法則}{こう|かん|ほう|そく}}\R{文字}{も|じ}は「どんな\R{数}{かず}でも」という\R{意味}{い|み}で\R{使}{つか}えるのです。ただし、ひき\R{算}{ざん}とわり\R{算}{ざん}では\R{入}{い}れかえられません（$5-2=3$ ですが、$2-5=-3$）。

もう1つ、\R{長方形}{ちょう|ほう|けい}の\R{面積}{めん|せき}で\R{考}{かんが}えましょう。\R{図}{ず}の\R{長方形}{ちょう|ほう|けい}\R{全体}{ぜん|たい}の\R{面積}{めん|せき}は、\R{何}{なん}\unit{cm^2}かな？

\noindent\begin{minipage}[c]{0.4\linewidth}\centering
\begin{tikzpicture}[x=1mm, y=1mm]
  \fill[bkBlue!20] (0,0) rectangle (14,28);
  \fill[bkOrange!25] (14,0) rectangle (35,28);
  \foreach \x in {7,14,21,28} \draw[bkLine, line width=0.4pt] (\x,0) -- (\x,28);
  \foreach \y in {7,14,21} \draw[bkLine, line width=0.4pt] (0,\y) -- (35,\y);
  \draw[bkInk, line width=0.9pt] (0,0) rectangle (35,28);
  \draw[bkInk, line width=0.9pt] (14,0) -- (14,28);
  \node[font=\sffamily\bfseries\small, text=bkBlue!80!black, fill=bkBlue!20, inner sep=1pt] at (7,14) {$2\times4$};
  \node[font=\sffamily\bfseries\small, text=bkOrange!80!black, fill=bkOrange!25, inner sep=1pt] at (24.5,14) {$3\times4$};
  \bkbrace{(0,29.5)}{(14,29.5)}{\qty{2}{cm}}{1}
  \bkbrace{(14,29.5)}{(35,29.5)}{\qty{3}{cm}}{1}
  \node[anchor=east, bk label] at (-1.5,14) {\qty{4}{cm}};
\end{tikzpicture}
\end{minipage}\hfill
\begin{minipage}[c]{0.57\linewidth}
\textbf{\R{考}{かんが}え\R{方}{かた}A}　\R{横}{よこ}の\R{長}{なが}さを\R{先}{さき}にたす（かっこの\R{中}{なか}が\R{先}{さき}）。\\
\hspace*{1\zw}$(2+3)\times4=5\times4=20$\par
\textbf{\R{考}{かんが}え\R{方}{かた}B}　2つの\R{長方形}{ちょう|ほう|けい}に\R{分}{わ}ける（×が\R{先}{さき}）。\\
\hspace*{1\zw}$2\times4+3\times4=8+12=20$\par
どちらも\qty{20}{cm^2}になりました。
\end{minipage}

\vspace{3pt}
ほかの\R{数}{かず}でも\R{確}{たし}かめましょう。\R{横}{よこ}が5と1、たてが3なら、$(5+1)\times3=18$、$5\times3+1\times3=15+3=18$。\R{横}{よこ}が1.5と2.5、たてが2なら、$(1.5+2.5)\times2=8$、$1.5\times2+2.5\times2=3+5=8$。いつでも\R{同}{おな}じになります。だから、\R{文字}{も|じ}を\R{使}{つか}って1\R{行}{ぎょう}で\R{書}{か}けます。
\[ (a+b)\times c=a\times c+b\times c \]
これを\textbf{\R{分配法則}{ぶん|ぱい|ほう|そく}}といいます。\bkidx{ぶんぱいほうそく}{\R{分配法則}{ぶん|ぱい|ほう|そく}}かっこの\R{外}{そと}の $c$ を、\R{中}{なか}の $a$ にも $b$ にも「\R{分}{わ}けて\R{配}{くば}る」から、この\R{名前}{な|まえ}がついています。ひき\R{算}{ざん}でも、$(a-b)\times c=a\times c-b\times c$ が\R{成}{な}り\R{立}{た}ちます。

\R{前}{まえ}のケーキとジュースも、\R{分配法則}{ぶん|ぱい|ほう|そく}で\R{確}{たし}かめられます。
\bkiilline{$(300+100)\times2=300\times2+100\times2=600+200=800$}
「ケーキ2こ\R{分}{ぶん}とジュース2\R{本}{ほん}\R{分}{ぶん}」と\R{考}{かんが}えても\R{同}{おな}じ800\R{円}{えん}です。かっこをつけた\R{式}{しき}と、つけずに×を2\R{回}{かい}\R{書}{か}いた\R{式}{しき}は、\R{同}{おな}じ\R{答}{こた}えになります。2-12では、このきまりで\R{文字式}{も|じ|しき}のかっこをはずします。

\begin{motto}[2]
\textbf{$x$ は、どこから\R{来}{き}たの？}\par
400\R{年}{ねん}ほど\R{前}{まえ}、フランスのデカルトは、わかっている\R{数}{かず}を $a$、$b$、$c$、わからない\R{数}{かず}を $x$、$y$、$z$ で\R{表}{あらわ}しました。この\R{書}{か}き\R{方}{かた}が\R{広}{ひろ}まって、\R{今}{いま}でも、わからない\R{数}{かず}には $x$ がよく\R{使}{つか}われます。
\end{motto}

\begin{taisetsu}
\R{文字}{も|じ}は、どんな\R{数}{かず}でも\R{入}{はい}る\R{箱}{はこ}。いつでも\R{使}{つか}える\R{式}{しき}が1\R{行}{ぎょう}で\R{書}{か}ける。\par\vspace{3pt}
\R{計算}{けい|さん}の\R{順番}{じゅん|ばん}　①かっこの\R{中}{なか} → ②×÷ → ③＋−（\R{同}{おな}じなかまは\R{左}{ひだり}から）。\R{累乗}{るい|じょう}は、かっこの\R{次}{つぎ}、×÷より\R{先}{さき}。\par\vspace{3pt}
×÷でできた\R{部分}{ぶ|ぶん}は1つの\R{量}{りょう}なので\R{先}{さき}に。＋−でできた\R{部分}{ぶ|ぶん}をひとまとまりにしたいときは、かっこでかこむ。
\end{taisetsu}

\begin{kotoba}
\small
\bkword{\R{文字式}{も|じ|しき}}{\R{文字}{も|じ}を\R{使}{つか}って\R{書}{か}いた\R{式}{しき}。$x=v\times t$ など。}
\bkword{\R{交換法則}{こう|かん|ほう|そく}}{$a+b=b+a$、$a\times b=b\times a$。\R{順番}{じゅん|ばん}を\R{入}{い}れかえても\R{答}{こた}えは\R{同}{おな}じ。}
\bkword{\R{分配法則}{ぶん|ぱい|ほう|そく}}{$(a+b)\times c=a\times c+b\times c$。かっこの\R{外}{そと}の\R{数}{かず}を、\R{中}{なか}のどちらにもかける。}
\end{kotoba}

\bkkigocards{\bkkigocard{$x$}{エックス}{\R{数}{かず}のかわりに\R{使}{つか}う\R{文字}{も|じ}。どんな\R{数}{かず}でも\R{入}{はい}る\R{箱}{はこ}。}{x}{2-11}}
{\bkkigocard{$(\ \ )$}{かっこ}{ひとまとまりにする\R{印}{しるし}。かっこの\R{中}{なか}を\R{先}{さき}に\R{計算}{けい|さん}する。}{parentheses}{2-11}}

\begin{column}{\R{電卓}{でん|たく}で、\R{計算}{けい|さん}の\R{順番}{じゅん|ばん}を\R{確}{たし}かめよう}
\leavevmode\label{bk:2-11:dentaku}\R{電卓}{でん|たく}で、\mbox{1＋3×100＝} と、ボタンを\R{順}{じゅん}におしてみましょう。301と\R{出}{で}る\R{電卓}{でん|たく}は、\R{計算}{けい|さん}の\R{順番}{じゅん|ばん}のきまりを\R{知}{し}っています。400と\R{出}{で}る\R{電卓}{でん|たく}は、おした\R{順}{じゅん}に\R{計算}{けい|さん}しています。どちらの\R{電卓}{でん|たく}でも、\R{先}{さき}に \mbox{3×100＝} で300を\R{出}{だ}してから1をたせば、きまりどおりの\R{答}{こた}えになります。2-18で\R{学}{まな}ぶ\R{表計算}{ひょう|けい|さん}や、2-19のプログラムは、きまりどおりに\R{計算}{けい|さん}します。
\end{column}

%% ---------------- 例題 ----------------
\clearpage
\bksection{reidai}{\R{例題}{れい|だい}でたしかめよう}

\begin{reidai}[きまりを\R{見}{み}つけて、\R{文字}{も|じ}の\R{式}{しき}を\R{作}{つく}る]
\R{自転車}{じ|てん|しゃ}が、スタートの\qty{10}{m}\R{先}{さき}から\qty{4}{m/s}で\R{進}{すす}みます。$t$ \R{秒後}{びょう|ご}の\R{位置}{い|ち} $x$ m を $t$ の\R{式}{しき}で\R{表}{あらわ}し、$t=5$ のときの\R{位置}{い|ち}を\R{求}{もと}めましょう。\par
\kangaekata $t=0$ で10、$t=1$ で14、$t=2$ で18。10に、1\R{秒}{びょう}ごとに4ずつふえる。$t$ \R{秒}{びょう}では、4が $t$ こ\R{分}{ぶん}ふえる。\par
\shiki \mbox{\R{位置}{い|ち}＝10＋4×\R{時間}{じ|かん}} だから、$x=10+4\times t$\par
\hspace*{2.6\zw}$t=5$ のとき、$x=10+4\times5=10+20=30$\quad（×を\R{先}{さき}に）\par
\kotae $x=10+4\times t$、\qty{30}{m}
\end{reidai}

\begin{jibun}
\qty{20}{L}の\R{水}{みず}が\R{入}{はい}った\R{水}{すい}そうに、1\R{分}{ぷん}に\qty{3}{L}ずつ\R{水}{みず}を\R{入}{い}れます。$t$ \R{分後}{ふん|ご}の\R{水}{みず}の\R{量}{りょう} $y$ L を $t$ の\R{式}{しき}で\R{表}{あらわ}し、$t=8$ のときを\R{求}{もと}めましょう。
\bkjans{$y=20+3\times t$、$t=8$ のとき $20+24=44$ で\qty{44}{L}}
\end{jibun}

\begin{jibun}
120ページの\R{本}{ほん}を、1\R{日}{にち}に8ページずつ\R{読}{よ}みます。$t$ \R{日後}{にち|ご}に\R{残}{のこ}っているページ\R{数}{すう}を $y$ として、$y$ を $t$ の\R{式}{しき}で\R{表}{あらわ}し、$t=5$ のときを\R{求}{もと}めましょう（\R{読}{よ}んだ\R{分}{ぶん}だけへるので、ひき\R{算}{ざん}になります）。
\bkjans{$y=120-8\times t$、$t=5$ のとき $120-40=80$ で80ページ}
\end{jibun}

\begin{reidai}[\R{計算}{けい|さん}の\R{順番}{じゅん|ばん}]
\R{計算}{けい|さん}しましょう。\quad（1）\nobreak $18-12\div3\times2$\quad（2）\nobreak $(8-2\times5)\times(-3)$\quad（3）\nobreak $20-7+5$\par
\kangaekata ①かっこの\R{中}{なか}　②×と÷　③＋と− の\R{順}{じゅん}に\R{計算}{けい|さん}します。\R{同}{おな}じなかまは\R{左}{ひだり}から。\R{先}{さき}に\R{計算}{けい|さん}する\R{所}{ところ}に\R{下線}{か|せん}を\R{引}{ひ}くと、まちがえにくくなります。\par
\shiki（1）÷と×は\R{左}{ひだり}から。$18-\underline{12\div3}\times2=18-\underline{4\times2}=18-8=10$\par
\hspace*{2.6\zw}（2）かっこの\R{中}{なか}も×が\R{先}{さき}。$(8-\underline{2\times5})\times(-3)=(8-10)\times(-3)$\par
\hspace*{4.6\zw}$=(-2)\times(-3)=6$\par
\hspace*{2.6\zw}（3）−と＋は\R{左}{ひだり}から。$\underline{20-7}+5=13+5=18$\par
\kotae（1）\nobreak 10\quad（2）\nobreak 6\quad（3）\nobreak 18
\end{reidai}

\begin{jibun}
\R{計算}{けい|さん}しましょう。\quad（1）\nobreak $20\div4\times2+6$\quad（2）\nobreak $(3-4\times2)\times(-2)$\quad（3）\nobreak $15-9+4$
\bkjans{（1）\nobreak 16\quad（2）\nobreak 10\quad（3）\nobreak 10}
\end{jibun}

\begin{reidai}[\R{文章}{ぶん|しょう}から1\R{行}{ぎょう}の\R{式}{しき}を\R{作}{つく}る：かっこは\R{要}{い}る？]
（1）1こ120\R{円}{えん}のパンを3こと、150\R{円}{えん}の\R{牛乳}{ぎゅう|にゅう}を1\R{本}{ぽん}\R{買}{か}って、1000\R{円}{えん}\R{出}{だ}しました。おつりはいくらですか。\par
（2）4\R{人}{にん}が、それぞれ300\R{円}{えん}のケーキと100\R{円}{えん}のジュースを\R{買}{か}います。\R{代金}{だい|きん}は\R{全部}{ぜん|ぶ}でいくらですか。\par
（3）\R{国語}{こく|ご}が80\R{点}{てん}、\R{算数}{さん|すう}が90\R{点}{てん}でした。2つの\R{平均}{へい|きん}は\R{何点}{なん|てん}ですか（\R{平均}{へい|きん}は、たしてから\R{個数}{こ|すう}でわる。2-6）。\par
\kangaekata ひとまとまりにしたい「\R{代金}{だい|きん}の\R{合計}{ごう|けい}」「1\R{人分}{にん|ぶん}」「\R{点数}{てん|すう}の\R{合計}{ごう|けい}」を、かっこでかこみます。\par
\shiki（1）\R{代金}{だい|きん}の\R{合計}{ごう|けい}を1000\R{円}{えん}からひく。$1000-(120\times3+150)=1000-(360+150)$\par
\hspace*{4.6\zw}$=1000-510=490$\par
\hspace*{2.6\zw}（2）1\R{人分}{にん|ぶん}が4つ\R{分}{ぶん}。$(300+100)\times4=400\times4=1600$\par
\hspace*{2.6\zw}（3）\R{合計}{ごう|けい}を2でわる。$(80+90)\div2=170\div2=85$\par
\kotae（1）490\R{円}{えん}\quad（2）1600\R{円}{えん}\quad（3）85\R{点}{てん}
\end{reidai}

\noindent（1）でかっこを\R{忘}{わす}れて $1000-120\times3+150$ と\R{書}{か}くと、$1000-360+150=790$ になります。これは\R{牛乳}{ぎゅう|にゅう}の150\R{円}{えん}を「はらう」のではなく「もらう」\R{式}{しき}です。

\begin{jibun}
（1）1\R{本}{ぽん}80\R{円}{えん}のえんぴつを5\R{本}{ほん}と、120\R{円}{えん}のノートを1\R{冊}{さつ}\R{買}{か}って、1000\R{円}{えん}\R{出}{だ}しました。おつりは？\quad（2）3\R{人}{にん}が、それぞれ250\R{円}{えん}のおにぎりと150\R{円}{えん}のお\R{茶}{ちゃ}を\R{買}{か}います。\R{代金}{だい|きん}は\R{全部}{ぜん|ぶ}で？\quad（3）\qty{38}{kg}と\qty{42}{kg}の\R{平均}{へい|きん}は？　どれも1\R{行}{ぎょう}の\R{式}{しき}で\R{書}{か}きましょう。
\bkjans{（1）\nobreak $1000-(80\times5+120)=480$ で480\R{円}{えん}\quad（2）\nobreak $(250+150)\times3=1200$ で1200\R{円}{えん}\quad（3）\nobreak $(38+42)\div2=40$ で\qty{40}{kg}}
\end{jibun}

\begin{reidai}[\R{式}{しき}が\R{表}{あらわ}す\R{場面}{ば|めん}を\R{読}{よ}む]
りんごは1こ120\R{円}{えん}、みかんは1こ50\R{円}{えん}です。（1）$120\times2+50\times3$、（2）$(120+50)\times3$ は、それぞれ\R{何}{なに}をいくつ\R{買}{か}ったときの\R{代金}{だい|きん}ですか。\R{代金}{だい|きん}も\R{求}{もと}めましょう。\par
\kangaekata ×でできた\R{部分}{ぶ|ぶん}は「1つ\R{分}{ぶん}×いくつ\R{分}{ぶん}」のまとまり。かっこは「ひとまとまり」の\R{印}{しるし}です。\par
\shiki（1）$120\times2$ はりんご2こ、$50\times3$ はみかん3こ。$240+150=390$\par
\hspace*{2.6\zw}（2）$120+50$ は、りんご1ことみかん1この\R{組}{くみ}。\par
\hspace*{4.6\zw}それが3\R{組}{くみ}なので、$170\times3=510$\par
\kotae（1）りんご2ことみかん3こで390\R{円}{えん}\quad（2）りんごとみかんを3こずつで510\R{円}{えん}
\end{reidai}

\begin{jibun}
ノートは1\R{冊}{さつ}150\R{円}{えん}、えんぴつは1\R{本}{ぽん}60\R{円}{えん}です。（1）$150\times4+60\times2$、（2）$(150+60)\times2$ は、それぞれ\R{何}{なに}をいくつ\R{買}{か}ったときの\R{代金}{だい|きん}ですか。\R{代金}{だい|きん}も\R{求}{もと}めましょう。
\bkjans{（1）\nobreak ノート4\R{冊}{さつ}とえんぴつ2\R{本}{ほん}で720\R{円}{えん}\quad（2）\nobreak ノートとえんぴつを2つずつで420\R{円}{えん}}
\end{jibun}

\begin{reidai}[\R{図}{ず}から、\R{文字}{も|じ}の\R{式}{しき}を\R{作}{つく}る]
\noindent\begin{minipage}[t]{0.71\linewidth}
つまようじで、\R{正三角形}{せい|さん|かく|けい}を\R{右}{みぎ}の\R{図}{ず}のようにつなげます。$n$ こつくるときの\R{本数}{ほん|すう}を $n$ の\R{式}{しき}で\R{表}{あらわ}し、$n=50$ のときを\R{求}{もと}めましょう。
\end{minipage}\hfill
\begin{minipage}[t]{0.27\linewidth}\vspace{0mm}\centering
\begin{tikzpicture}[x=0.8mm, y=0.8mm]
  \bkiilpick{(0,0)}{(6,10.4)}{bkOrange}
  \foreach \k in {0,1,2} {\pgfmathsetmacro\xa{\k*12}
    \bkiilpick{(\xa,0)}{(\xa+12,0)}{bkBlue!70}\bkiilpick{(\xa+12,0)}{(\xa+6,10.4)}{bkBlue!70}}
  \foreach \k in {0,1} {\pgfmathsetmacro\xa{\k*12}
    \bkiilpick{(\xa+6,10.4)}{(\xa+18,10.4)}{bkBlue!70}\bkiilpick{(\xa+18,10.4)}{(\xa+12,0)}{bkBlue!70}}
  \node[font=\sffamily\footnotesize, anchor=north] at (18,-1.5) {$n=5$};
\end{tikzpicture}
\end{minipage}\par
\kangaekata \R{左}{ひだり}はしの1\R{本}{ぽん}（オレンジ）に、\R{三角形}{さん|かく|けい}1こにつき2\R{本}{ほん}ずつ（\R{青}{あお}）がつながる。\par
\shiki \R{本数}{ほん|すう}＝$1+2\times n$。$n=50$ のとき、$1+2\times50=1+100=101$\quad（×を\R{先}{さき}に）\par
\kotae $1+2\times n$（\R{本}{ほん}）、101\R{本}{ぽん}
\end{reidai}

\begin{jibun}
\noindent\begin{minipage}[t]{0.71\linewidth}
\R{正六角形}{せい|ろっ|かく|けい}を\R{右}{みぎ}の\R{図}{ず}のようにつなげます。$n$ このときの\R{本数}{ほん|すう}を $n$ の\R{式}{しき}で\R{表}{あらわ}し、$n=30$ のときを\R{求}{もと}めよう。
\end{minipage}\hfill
\begin{minipage}[t]{0.27\linewidth}\vspace{0mm}\centering
\begin{tikzpicture}[x=0.75mm, y=0.75mm]
  \foreach \k in {0,1,2} {\begin{scope}[shift={(\k*10.4,0)}]
    \foreach \a in {90,150,210,270,330,30} {\pgfmathsetmacro\b{\a+60}\bkiilpick{(\a:6)}{(\b:6)}{bkiilwood}}
  \end{scope}}
  \node[font=\sffamily\footnotesize, anchor=north] at (10.4,-6.3) {$n=3$};
\end{tikzpicture}
\end{minipage}
\bkjans{$1+5\times n$（\R{本}{ほん}）、$n=30$ のとき $1+150=151$ で151\R{本}{ぽん}}
\end{jibun}

\begin{motto}[2]
\textbf{きまりを\R{見}{み}つけて\R{式}{しき}を\R{作}{つく}る3つの\R{手順}{て|じゅん}}\par
つまようじの\R{問題}{もん|だい}は、いつも\R{同}{おな}じ\R{手順}{て|じゅん}で\R{式}{しき}にできます。\R{例題}{れい|だい}の\R{正三角形}{せい|さん|かく|けい}で\R{確}{たし}かめましょう。
\begin{enumerate}[label={\bkcircled[bkGray]{\arabic*}}, leftmargin=2.4\zw, itemsep=1pt, topsep=2pt]
\item $n=1$、2、3、…のときの\R{本数}{ほん|すう}を、\R{表}{ひょう}に\R{書}{か}く。
\item 1こふえるたびに、\R{何本}{なん|ぼん}ふえるか（ふえる\R{数}{かず}）を\R{見}{み}る。
\item $n=1$ の\R{本数}{ほん|すう}から、ふえる\R{数}{かず}を1\R{回}{かい}ひいて、「はじめの\R{数}{かず}」を\R{見}{み}つける。
\end{enumerate}
\begin{bkcenter}
\begin{bktable}{colspec={Q[l,m,wd=36mm] *{5}{Q[c,m,wd=10mm]}}, rowsep=1pt, colsep=4pt, cells={font=\small}, row{1}={bg=bkGray!12}}
\R{三角形}{さん|かく|けい}の\R{数}{かず} $n$ & 1 & 2 & 3 & 4 & 5 \\
\R{本数}{ほん|すう} & 3 & 5 & 7 & 9 & 11 \\
\end{bktable}
\end{bkcenter}
\R{表}{ひょう}を\R{右}{みぎ}へ1つ\R{進}{すす}むたびに2\R{本}{ほん}ふえるので、ふえる\R{数}{かず}は2。はじめの\R{数}{かず}は $3-2=1$。だから\R{式}{しき}は「はじめの\R{数}{かず}＋ふえる\R{数}{かず}×$n$」で、$1+2\times n$ です。$n=5$ を\R{入}{い}れると $1+10=11$ で、\R{表}{ひょう}と\R{合}{あ}います。\R{例題}{れい|だい}1の $x=10+4\times t$ も、はじめの\R{数}{かず}10に、ふえる\R{数}{かず}4が $t$ こ\R{分}{ぶん}たされた、\R{同}{おな}じ\R{形}{かたち}の\R{式}{しき}です。
\end{motto}

\begin{reidai}[かっこをつける\R{場所}{ば|しょ}で、\R{答}{こた}えが\R{変}{か}わる]
$6+4\times3-2$ に、かっこを1\R{組}{くみ}だけつけます。\R{答}{こた}えがいちばん\R{大}{おお}きくなるのは、どこにつけたときですか。\par
\kangaekata かっこをつけられる\R{場所}{ば|しょ}を\R{全部}{ぜん|ぶ}ためして、くらべます。かっこがないときも\R{計算}{けい|さん}しておきます。\par
\shiki かっこなし：$6+12-2=16$\par
\hspace*{2.6\zw}$(6+4)\times3-2=10\times3-2=30-2=28$\par
\hspace*{2.6\zw}$6+4\times(3-2)=6+4\times1=10$\par
\hspace*{2.6\zw}$(6+4\times3)-2=18-2=16$、\ $6+(4\times3-2)=6+10=16$\par
\kotae $(6+4)\times3-2=28$
\end{reidai}

\begin{jibun}
$1+2\times5-3$ に、かっこを1\R{組}{くみ}だけつけて、\R{答}{こた}えをできるだけ\R{大}{おお}きくしましょう。そのときの\R{答}{こた}えも\R{書}{か}きましょう。
\bkjans{$(1+2)\times5-3=12$}
\end{jibun}

\begin{reidai}[\R{分配法則}{ぶん|ぱい|ほう|そく}で、くふうして\R{計算}{けい|さん}する]
くふうして\R{計算}{けい|さん}しましょう。\quad（1）\nobreak $102\times6$\quad（2）\nobreak $99\times8$\par
\kangaekata 102を $100+2$、99を $100-1$ と\R{分}{わ}けて、\R{分配法則}{ぶん|ぱい|ほう|そく}を\R{使}{つか}います。100のかけ\R{算}{ざん}は、\R{暗算}{あん|ざん}でかんたんにできるからです。\par
\shiki（1）$102\times6=(100+2)\times6=100\times6+2\times6=600+12=612$\par
\hspace*{2.6\zw}（2）$99\times8=(100-1)\times8=100\times8-1\times8=800-8=792$\par
\kotae（1）\nobreak 612\quad（2）\nobreak 792
\end{reidai}

\begin{jibun}
くふうして\R{計算}{けい|さん}しましょう。\quad（1）\nobreak $103\times7$\quad（2）\nobreak $98\times5$
\bkjans{（1）\nobreak $(100+3)\times7=700+21=721$\quad（2）\nobreak $(100-2)\times5=500-10=490$}
\end{jibun}

%% ---------------- まちがえやすいところ ----------------
\bksection{tsumazuki}{まちがえやすいところをチェック}

\begin{tsumazuki}{\R{左}{ひだり}から\R{順}{じゅん}に\R{計算}{けい|さん}してしまう}
\bkng{$6+4\times2=10\times2=20$}
\bkok{$6+4\times2=6+8=14$}
×と÷が\R{先}{さき}。おした\R{順}{じゅん}に\R{計算}{けい|さん}する\R{電卓}{でん|たく}に\R{注意}{ちゅう|い}（\pageref{bk:2-11:dentaku}ページ）。
\end{tsumazuki}

\begin{tsumazuki}{×を、÷より\R{先}{さき}に\R{計算}{けい|さん}してしまう}
\bkng{$24\div4\times2=24\div8=3$}
\bkok{$24\div4\times2=6\times2=12$}
×と÷は\R{同}{おな}じなかま。×と÷だけなら、\R{左}{ひだり}から\R{順}{じゅん}に\R{計算}{けい|さん}します。
\end{tsumazuki}

\begin{tsumazuki}{＋を、−より\R{先}{さき}に\R{計算}{けい|さん}してしまう}
\bkng{$12-5+3=12-8=4$}
\bkok{$12-5+3=7+3=10$}
−5の−は、5だけについた\R{印}{しるし}。＋と−だけなら、\R{左}{ひだり}から\R{順}{じゅん}に\R{計算}{けい|さん}します。
\end{tsumazuki}

\begin{tsumazuki}{9\R{秒}{びょう}と10\R{秒}{びょう}の\R{平均}{へい|きん}で、かっこを\R{忘}{わす}れる}
\bkng{$9+10\div2=9+5=14$}
\bkok{$(9+10)\div2=19\div2=9.5$}
\R{平均}{へい|きん}は、たしてからわる（2-6）。9\R{秒}{びょう}と10\R{秒}{びょう}の\R{平均}{へい|きん}が14\R{秒}{びょう}になるのは、おかしいとわかります。
\end{tsumazuki}

\begin{tsumazuki}{かっこの\R{中}{なか}より\R{先}{さき}に、かっこの\R{外}{そと}を\R{計算}{けい|さん}してしまう}
\bkng{$(2+3)\times5=2+15=17$}
\bkok{$(2+3)\times5=5\times5=25$}
かっこの\R{中}{なか}は1つの\R{数}{かず}のかわり。いちばん\R{先}{さき}に\R{計算}{けい|さん}します。
\end{tsumazuki}



%% ---------------- 物理・工学 ----------------
\begin{butsuri}
\leavevmode\label{bk:2-11:butsuri}\textbf{\R{車}{くるま}が\R{止}{と}まるまでの\R{距離}{きょ|り}}\quad \R{運転手}{うん|てん|しゅ}が\R{危険}{き|けん}に\R{気}{き}づいてから、ブレーキをふむまでに、およそ0.75\R{秒}{びょう}かかります。そのあいだも、\R{車}{くるま}はそれまでの\R{速}{はや}さで\R{進}{すす}みます。この\R{距離}{きょ|り}を\textbf{\R{空走距離}{くう|そう|きょ|り}}といいます。\R{速}{はや}さを $v$（m/s）、ブレーキをふむまでの\R{時間}{じ|かん}を $t$（s）とすると、\R{空走距離}{くう|そう|きょ|り} $x$（m）は \mbox{$x=v\times t$} で\R{計算}{けい|さん}できます。$t=0.75$ として、いろいろな\R{速}{はや}さを\R{入}{い}れてみましょう。
\begin{bkfigure}
\begin{tikzpicture}[x=1mm, y=1mm]
  \fill[bkGray!25] (0,0) rectangle (118,9);
  \foreach \x in {6,22,...,102} \fill[white] (\x,4.1) rectangle ++(8,0.8);
  \bkcar{(4,9)}{0.62}{bkBlue}
  \fill[bkRed] (113,13) circle (3.2); \fill[white] (112,14) circle (0.9);
  \foreach \xx/\t in {27/\R{気}{き}づく, 57/ブレーキをふむ, 104/\R{止}{と}まる} {
    \draw[bkInk, line width=0.7pt, dashed] (\xx,9) -- (\xx,25);
    \node[font=\sffamily\footnotesize, anchor=south] at (\xx,25.5) {\t};}
  \draw[bk arrow, bkOrange, line width=1.1pt, {Stealth[length=5pt]}-{Stealth[length=5pt]}] (27.6,15) -- node[above=0.5mm, font=\sffamily\bfseries\footnotesize, text=bkOrange!85!black] {\R{空走距離}{くう|そう|きょ|り}} (56.4,15);
  \draw[bk arrow, bkPurple, line width=1.1pt, {Stealth[length=5pt]}-{Stealth[length=5pt]}] (57.6,15) -- node[above=0.5mm, font=\sffamily\bfseries\footnotesize, text=bkPurple!85!black] {\R{制動距離}{せい|どう|きょ|り}} (103.4,15);
\end{tikzpicture}
\bkfigcaption{\R{制動距離}{せい|どう|きょ|り}は、ブレーキがきいてから\R{止}{と}まるまでの\R{距離}{きょ|り}。}
\end{bkfigure}
\begin{bkcenter}
\begin{bktable}{colspec={Q[l,m,wd=28mm] *{3}{Q[c,m,wd=27mm]}}, rowsep=1pt, cells={font=\small}, row{1}={bg=bkPurple!12, font=\footnotesize}}
\R{速}{はや}さ $v$ & {\qty{10}{m/s}\\（\R{時速}{じ|そく}\qty{36}{km}）} & {\qty{20}{m/s}\\（\R{時速}{じ|そく}\qty{72}{km}）} & {\qty{30}{m/s}\\（\R{時速}{じ|そく}\qty{108}{km}）} \\
\R{空走距離}{くう|そう|きょ|り} $v\times0.75$ & \qty{7.5}{m} & \qty{15}{m} & \qty{22.5}{m} \\
\end{bktable}
\end{bkcenter}
\R{時速}{じ|そく}は、2-10の\R{単位}{たん|い}つきの1で\R{換算}{かん|さん}しました。たとえば \mbox{\qty{10}{m/s}×\qty{3600}{s/h}÷\qty{1000}{m/km}＝\qty{36}{km/h}}（sとmが\R{約分}{やく|ぶん}されて\R{消}{き}える）。\R{速}{はや}さが2\R{倍}{ばい}になると、\R{空走距離}{くう|そう|きょ|り}も2\R{倍}{ばい}。\R{自動}{じ|どう}ブレーキのある\R{車}{くるま}のコンピュータも、はかった\R{距離}{きょ|り}や\R{速}{はや}さを\R{文字}{も|じ}の\R{式}{しき}に\R{入}{い}れて、ぶつからないかを\R{計算}{けい|さん}しつづけています。
\end{butsuri}



%% ---------------- 章末問題 ----------------
\clearpage
\bkendproblems
\bkproblemgroup{\R{基本}{き|ほん}}{1}

\bkproblem \R{次}{つぎ}の\R{計算}{けい|さん}をしましょう。\par
（1）\nobreak $8+6\times3$\quad（2）\nobreak $(8+6)\times3$\quad（3）\nobreak $30-18\div6$\quad（4）\nobreak $16\div4\times2$\par
（5）\nobreak $(15-3)\div(2+4)$\quad（6）\nobreak $20-8+3$\quad（7）\nobreak $2+3\times5$ と $(2+3)\times5$
\bkpans{（1）\nobreak 26\quad（2）\nobreak 42\quad（3）\nobreak 27\quad（4）\nobreak 8\quad（5）\nobreak 2\quad（6）\nobreak 15\quad（7）\nobreak 17と25}

\bkproblem \R{次}{つぎ}の\R{計算}{けい|さん}をしましょう。\par
（1）\nobreak $-5+2\times(-3)$\quad（2）\nobreak $(-4)\times(6-9)$\quad（3）\nobreak $10-(-2)\times(-4)$
\bkpans{（1）\nobreak $-11$\quad（2）\nobreak 12\quad（3）\nobreak 2}

\bkproblem 500\R{円}{えん}\R{出}{だ}して、1こ80\R{円}{えん}のガムを2こと、60\R{円}{えん}のあめを1こ\R{買}{か}いました。おつりを1\R{行}{ぎょう}の\R{式}{しき}で\R{書}{か}いて\R{求}{もと}めましょう。
\bkpans{$500-(80\times2+60)=500-220=280$ で280\R{円}{えん}}

\bkproblem \R{長}{なが}さ\qty{30}{cm}のろうそくが、\R{火}{ひ}をつけると1\R{分}{ぷん}に\qty{2}{cm}ずつ\R{短}{みじか}くなります。$t$ \R{分後}{ふん|ご}のろうそくの\R{長}{なが}さを $y$ cm として、$y$ を $t$ の\R{式}{しき}で\R{表}{あらわ}しましょう。また、$t=10$ のときの\R{長}{なが}さを\R{求}{もと}めましょう。
\bkpans{$y=30-2\times t$、$t=10$ のとき $30-20=10$ で\qty{10}{cm}}

\bkproblem \noindent\begin{minipage}[t]{0.58\linewidth}
つまようじで、\R{正方形}{せい|ほう|けい}を\R{右}{みぎ}の\R{図}{ず}のように2\R{段}{だん}にして\R{横}{よこ}につなげます。たてに2こならんだ\R{正方形}{せい|ほう|けい}の\R{組}{くみ}を $n$ \R{組}{くみ}つくるときの\R{本数}{ほん|すう}を、$n$ の\R{式}{しき}で\R{表}{あらわ}しましょう。また、$n=10$ のときの\R{本数}{ほん|すう}を\R{求}{もと}めましょう。
\end{minipage}\hfill
\begin{minipage}[t]{0.34\linewidth}\vspace{0mm}\centering
\begin{tikzpicture}[x=1mm, y=1mm]
  \bkiilsq{(0,0)}{3}{9}{bkiilwood}{bkiilwood}
  \foreach \k in {1,2,3} {\bkiilpick{({(\k-1)*9},18)}{({\k*9},18)}{bkiilwood}\bkiilpick{({\k*9},9)}{({\k*9},18)}{bkiilwood}}
  \bkiilpick{(0,9)}{(0,18)}{bkiilwood}
  \node[font=\sffamily\footnotesize, anchor=north] at (13.5,-1.5) {$n=3$};
\end{tikzpicture}
\end{minipage}
\bkpans{$2+5\times n$（\R{本}{ほん}）、$n=10$ のとき52\R{本}{ほん}}

\bkproblem \R{分配法則}{ぶん|ぱい|ほう|そく}を\R{使}{つか}って、くふうして\R{計算}{けい|さん}しましょう。\quad（1）\nobreak $101\times9$\quad（2）\nobreak $49\times4$
\bkpans{（1）\nobreak $(100+1)\times9=909$\quad（2）\nobreak $(50-1)\times4=196$}

\bkproblem \R{貸}{か}し\R{自転車}{じ|てん|しゃ}の\R{料金}{りょう|きん}は、200\R{円}{えん}に、1\R{時間}{じ|かん}ごとに100\R{円}{えん}ずつたしていきます。$t$ \R{時間}{じ|かん}\R{使}{つか}ったときの\R{料金}{りょう|きん}を $y$ \R{円}{えん}として、$y$ を $t$ の\R{式}{しき}で\R{表}{あらわ}しましょう。また、$t=3$ のときの\R{料金}{りょう|きん}を\R{求}{もと}めましょう。
\bkpans{$y=200+100\times t$、$t=3$ のとき500\R{円}{えん}}

\bkproblem 1こ60\R{円}{えん}のみかんを5こ\R{買}{か}って、500\R{円}{えん}\R{出}{だ}しました。おつりを1\R{行}{ぎょう}の\R{式}{しき}で\R{書}{か}いて\R{求}{もと}めましょう。かっこは\R{要}{い}りますか。
\bkpans{$500-60\times5=500-300=200$ で200\R{円}{えん}。ひくのは $60\times5$ という1つのまとまりで、×が\R{先}{さき}なので、かっこは\R{要}{い}らない}

\bkproblem ジュースは1\R{本}{ぽん}120\R{円}{えん}、パンは1こ150\R{円}{えん}です。（1）$120\times3+150\times2$、（2）$(120+150)\times2$ は、それぞれ\R{何}{なに}をいくつ\R{買}{か}ったときの\R{代金}{だい|きん}ですか。\R{代金}{だい|きん}も\R{求}{もと}めましょう。
\bkpans{（1）\nobreak ジュース3\R{本}{ぼん}とパン2こで660\R{円}{えん}\quad（2）\nobreak ジュースとパンを2つずつで540\R{円}{えん}}

\bkproblem \R{算数}{さん|すう}のテストが70\R{点}{てん}、80\R{点}{てん}、90\R{点}{てん}でした。\R{平均}{へい|きん}を1\R{行}{ぎょう}の\R{式}{しき}で\R{書}{か}いて\R{求}{もと}めましょう。
\bkpans{$(70+80+90)\div3=240\div3=80$ で80\R{点}{てん}}


\bkproblemgroup{チャレンジ}{2}

\bkproblem \R{車}{くるま}が\R{止}{と}まるまでの\R{空走距離}{くう|そう|きょ|り}は $x=v\times t$ で\R{計算}{けい|さん}できます（\pageref{bk:2-11:butsuri}ページ）。ブレーキをふむまでの\R{時間}{じ|かん}を $t=0.75$（\R{秒}{びょう}）として、\R{速}{はや}さが\qty{12}{m/s}のときと\qty{25}{m/s}のときの\R{空走距離}{くう|そう|きょ|り}を\R{求}{もと}めましょう。
\bkpans{\qty{12}{m/s}のとき\qty{9}{m}、\qty{25}{m/s}のとき\qty{18.75}{m}}

\bkproblem $2+3\times4-1$ に、かっこを1\R{組}{くみ}だけつけて、\R{答}{こた}えをできるだけ\R{大}{おお}きくしましょう。そのときの\R{答}{こた}えも\R{書}{か}きましょう。
\bkpans{$(2+3)\times4-1=19$}

\bkproblem 3\R{回}{かい}はかった50 m\R{走}{そう}のタイムは、9.4\R{秒}{びょう}、9.8\R{秒}{びょう}、9.6\R{秒}{びょう}でした。\R{平均}{へい|きん}を1\R{行}{ぎょう}の\R{式}{しき}で\R{書}{か}いて、\R{計算}{けい|さん}しましょう。
\begin{bkcenter}
\begin{tikzpicture}[x=1mm, y=1mm]
  \foreach \t/\xx/\k in {9.4/0/1\R{回目}{かい|め}, 9.8/34/2\R{回目}{かい|め}, 9.6/68/3\R{回目}{かい|め}} {\bkdisplay{(\xx,0)}{26}{\t}\node[bk label, anchor=north] at (\xx+13,-1.5) {\k};}
\end{tikzpicture}
\end{bkcenter}
\bkpans{$(9.4+9.8+9.6)\div3=28.8\div3=9.6$ で9.6\R{秒}{びょう}}

\bkproblem 1こ\qty{4}{kg}の\R{箱}{はこ}Aと、1こ\qty{2}{kg}の\R{箱}{はこ}Bがあります。（1）$4\times3+2\times5$、（2）$(4+2)\times5$ は、それぞれ\R{何}{なに}をいくつ\R{運}{はこ}んだときの\R{重}{おも}さですか。\R{重}{おも}さも\R{求}{もと}めましょう。
\bkpans{（1）\nobreak Aを3ことBを5こで\qty{22}{kg}\quad（2）\nobreak AとBを1こずつ\R{組}{くみ}にして5\R{組}{くみ}で\qty{30}{kg}}

\bkproblem $8-2\times3+1$ に、かっこを1\R{組}{くみ}だけつけて、\R{答}{こた}えをできるだけ\R{小}{ちい}さくしましょう。そのときの\R{答}{こた}えも\R{書}{か}きましょう。
\bkpans{$8-2\times(3+1)=8-8=0$}

\begin{column}{\R{累乗}{るい|じょう}があるときは？　\R{世界}{せ|かい}の\R{覚}{おぼ}え\R{方}{かた}}
2-7の\R{累乗}{るい|じょう}がある\R{式}{しき}では、\R{累乗}{るい|じょう}をかっこの\R{次}{つぎ}、×÷より\R{先}{さき}に\R{計算}{けい|さん}します。$3^2$ は「3を2\R{回}{かい}かけた1つのまとまり」だからです。たとえば $2\times3^2=2\times9=18$ で、$6^2=36$ ではありません。アメリカでは、この\R{順番}{じゅん|ばん}を、かっこ（Parentheses）・\R{累乗}{るい|じょう}（Exponents）・×÷（Multiplication, Division）・＋−（Addition, Subtraction）の\R{頭文字}{かしら|も|じ}をとって「PEMDAS（ペムダス）」と\R{覚}{おぼ}えます。\R{国}{くに}がちがっても、\R{計算}{けい|さん}の\R{順番}{じゅん|ばん}は\R{同}{おな}じです。
\end{column}

\begin{furikaeri}
\begin{checklist}
\item ことばの\R{式}{しき}を、□や\R{文字}{も|じ}を\R{使}{つか}った\R{式}{しき}で\R{表}{あらわ}せる。\hfill{\footnotesize\sffamily\color{bkTheme}→\pageref{bk:2-11:moji}ページ}
\item \R{計算}{けい|さん}の\R{順番}{じゅん|ばん}のきまりと、そのわけがわかる。\hfill{\footnotesize\sffamily\color{bkTheme}→\pageref{bk:2-11:junban}ページ}
\item どんな\R{数}{かず}でも\R{成}{な}り\R{立}{た}つきまりを、\R{文字}{も|じ}で1\R{行}{ぎょう}に\R{書}{か}ける。\hfill{\footnotesize\sffamily\color{bkTheme}→\pageref{bk:2-11:kimari}ページ}
\end{checklist}
\end{furikaeri}

\bkfukidashi{sora}{\R{計算}{けい|さん}の\R{順番}{じゅん|ばん}は、\R{式}{しき}の\R{意味}{い|み}どおりに\R{計算}{けい|さん}するためのやくそくなんだね。\R{次}{つぎ}の2-12では、\R{文字式}{も|じ|しき}をもっと\R{短}{みじか}く\R{書}{か}くやくそくを\R{学}{まな}ぶよ。}
